% ------------------------------------------------------------------
% Fig. 2 -- How the exhaustive reference front was built.
% ------------------------------------------------------------------
\begin{figure}[t]
\centering
\begin{tikzpicture}[node distance=4.5mm]

\node[parambox, minimum width=64mm] (space)
     {\textbf{The four settings that matter}\hspace{1mm}
      $L$, $w$, $b$, $\delta$\\[1pt]
      $\tau$ and $\pi$ left out --- they change only latency};

\node[box, below=of space, minimum width=64mm] (grid)
     {\textbf{Try every point on a dense grid}\\[1pt]
      $3 \times 151 \times 4 \times 51 = 92{,}412$ designs};

\node[modelbox, below=of grid, minimum width=64mm] (sched)
     {\textbf{Give each one its best possible latency}\\[1pt]
      $207{,}927$ scheduler simulations, $248$\,s in total};

\node[box, below=of sched, minimum width=64mm] (filter)
     {\textbf{Discard anything slower than 100\,ms}\\[1pt]
      $26{,}820$ left ($29.0\%$)};

\node[box, below=of filter, minimum width=64mm] (nds)
     {\textbf{Keep only the best trade-offs}\\[1pt]
      designs no other design beats on both goals};

\node[metricbox, below=of nds, minimum width=64mm] (front)
     {\textbf{Reference set: $108$ designs}\\[1pt]
      score $0.016821$};

\node[box, fill=yellow!10, draw=black!55, below=of front, minimum width=64mm] (use)
     {This is the ceiling every method is measured against};

\draw[flow] (space)  -- (grid);
\draw[flow] (grid)   -- (sched);
\draw[flow] (sched)  -- (filter);
\draw[flow] (filter) -- (nds);
\draw[flow] (nds)    -- (front);
\draw[flow] (front)  -- (use);

\end{tikzpicture}
\caption{Building the reference set. Because the two scheduling settings
$\tau$ and $\pi$ affect nothing but latency, each combination of the other
four can be given its best achievable latency and the rest of the space
enumerated exactly. The resulting $108$ designs are the true best trade-offs
available, which lets us report how close a method got to the best possible
answer rather than only how it compared to another method.}
\label{fig:reffront}
\end{figure}
